% !TEX program = xelatex
\documentclass{英语试卷}
\begin{document}
\examtitle{2025年合肥市高三第二次教学质量检测}{英语}
\begin{notice}
\par 1. 答卷前，考生务必将自己的姓名、准考证号填写在答题卡上，\par
\par 2. 回答选择题时，选出每小题答案后，用铅笔把答题卡上对应题目的答案标号涂黑。如需改动，用橡皮擦干净后，再选涂其他答案标号。回答非选择题时，将答案写在答题卡上，写在本试卷上无效。\par
\par 3. 考试结束后，将本试卷和答题卡一并交回。\par
\end{notice}
\parttitle{第一部分 听力（共两节，满分 30 分）}
\par 做题时，先将答案标在试卷上。录音内容结束后，你将有两分钟的时间将试卷上的答案转涂到答题卡上。\par
\sectitle{第一节（共5小题；每小题1.5分，满分7.5分）}
\par 听下面 5 段对话。每段对话后有一个小题，从题中所给的 A、B、C 三个选项中选出最佳选项，并标在试卷的相应位置。听完每段对话后，你都有 10 秒钟的时间来回答有关小题和阅读下一小题。每段对话仅读一遍。\par
\question{例：How much is the shirt?}
\begin{choices}(3)
  \item £ 19.15.
  \item £ 9.18.
  \item £ 9.15. 答案是 C,
\end{choices}
\question{1. What has Sarah done recently?}
\begin{choices}(3)
  \item She has changed her job.
  \item She has taken a course.
  \item She has started working out.
\end{choices}
\question{2. What are the speakers mainly talking about?}
\begin{choices}(3)
  \item A weekend plan
  \item A neighborhood event.
  \item A street performance.
\end{choices}
\question{3. What time does the woman plan to arrive tomorrow?}
\begin{choices}(3)
  \item At 8:50.
  \item At 9:00.
  \item At 9:10.
\end{choices}
\question{4. What is the probable relationship between James and Judy?}
\begin{choices}(3)
  \item Husband and wife.
  \item Fellow workers.
  \item Classmates.
\end{choices}
\question{5. How does the man find Lena's work?}
\begin{choices}(3)
  \item It's satisfactory.
  \item It's average.
  \item It's disappointing.
\end{choices}
\sectitle{第二节（共 15 小题；每小题 1.5 分，满分 22.5 分）}
\par 听下面 5 段对话或独白。每段对话或独白后有几个小题，从题中所给的 A、B、C 三个选项中选出最佳选项，并标在试卷的相应位置。听每段对话或独白前，你将有时间阅读各个小题，每小题 5 秒钟；听完后，各小题将给出 5 秒钟的作答时间。每段对话或独白读两遍。\par
\par 听第 6 段材料, 回答第 6、7 题。\par
\question{6. What does the man remind the woman of?}
\begin{choices}(3)
  \item The date of the field trip.
  \item The confirmation page.
  \item The submission deadline.
\end{choices}
\question{7. What aspect of the field trip particularly interests the woman?}
\begin{choices}(3)
  \item The lectures.
  \item The exhibitions.
  \item The discussions.
\end{choices}
\par 听第 7 段材料，回答第 8 至 10 题。\par
\question{8. Where did the man go last year?}
\begin{choices}(3)
  \item To the Eiffel Tower in Paris.
  \item To the Palace Museum in Beijing.
  \item To the Colosseum in Rome.
\end{choices}
\question{9. What does the woman say about her experience in Rome last summer?}
\begin{choices}(3)
  \item It is eye-opening.
  \item It is stress-relieving.
  \item It is heart-warming.
\end{choices}
\question{10. What do the speakers agree on?}
\begin{choices}(3)
  \item The beauty of the Eiffel Tower.
  \item The brilliance of ancient Roman builders.
  \item The importance of visiting historic places.
\end{choices}
\par 听第 8 段材料, 回答第 11 至 13 题,\par
\question{11. Why does Finland return to paper textbooks according to the woman?}
\begin{choices}(3)
  \item To improve kids' comprehension.
  \item To limit kids' time on digital devices.
  \item To offer kids' hands-on experience.
\end{choices}
\question{12. How does Sam feel about paper textbooks?}
\begin{choices}(3)
  \item They harm the environment.
  \item They narrow the digital gap.
  \item They help students concentrate.
\end{choices}
\question{13. What will the speakers suggest to the school board?}
\begin{choices}(3)
  \item Reducing paper consumption.
  \item Switching to paper textbooks.
  \item Using a mixed model of textbooks.
\end{choices}
\par 听第 9 段材料, 回答第 14 至 17 题,\par
\question{14. What is the man doing?}
\begin{choices}(3)
  \item Attending a lecture.
  \item Hosting a workshop.
  \item Conducting an interview.
\end{choices}
\question{15. What may keep Ruby productive?}
\begin{choices}(3)
  \item Her passion for her work.
  \item Her gift for time management.
  \item Her desire to realize her dream.
\end{choices}
\question{16. Which does Ruby think helps her look for inspiration?}
\begin{choices}(3)
  \item Walking in a park.
  \item Playing light music.
  \item Talking with friends.
\end{choices}
\question{17. What is Ruby's primary advice for young readers?}
\begin{choices}(3)
  \item Write with passion.
  \item Be patient with writing.
  \item Learn from other writers.
\end{choices}
\par 听第 10 段材料，回答第 18 至 20 题，\par
\question{18. What is Replika?}
\begin{choices}(3)
  \item A digital companion.
  \item A chat-based website.
  \item A social networking app.
\end{choices}
\question{19. How should AI success be measured according to the speaker?}
\begin{choices}(3)
  \item By its emotional diversity.
  \item By its impact on users' emotions.
  \item By screen time and user engagement.
\end{choices}
\question{20. Why does the speaker give the talk?}
\begin{choices}(2)
  \item To introduce the latest technology behind AI. 3. To explore the role of AI in social connections.
  \item To discuss the benefits and risks of AI friends.
\end{choices}
\parttitle{第二部分 阅读理解（共两节，满分 50 分）}
\sectitle{第一节（共15小题；每小题2.5分，满分37.5分）}
\par 阅读下列短文，从每题所给的 A、B、C 和 D 四个选项中，选出最佳答案。\par
\par The Splendid Colors: Chinese Costume Culture Exhibition\par
\par Date: Oct. 1-Dec.2 Location: South Exhibition Hall 12, 6/F, the China National Arts and Crafts Museum in Beijing Highlights\par
\par This exhibition showcases traditional costumes from more than 30 ethnic (民族的) groups across various dynasties, emphasizing the unity within Chinese culture. Visitors can gain insights into their daily activities and infer the environmental conditions based on the materials and techniques used in the costumes. This is why experts nowadays refer to these costumes as "wearable history books".\par
\par Online Reservation\par
\par All visitors can make a reservation 1-7 days in advance through our official website or official WeChat platform. All visitors shall arrive at the museum on time according to the reserved time schedule. If you miss the appointment, entry will be denied. If you cannot visit on time, please cancel at least 2 hours beforehand. Late cancellations will not get your money back.\par
\par We do not currently accept group reservations.\par
\par Entry Instructions\par
\par Upon entering the museum, please present the original ID card used for making the reservation and the generated QR code(二维码). International visitors must present their passport along with the QR code. All visitors entering the museum shall go through a security screening. Visitors with access requirements and their accompanying family members or personal carers can skip the queues to enter the museum.\par
\question{21. What is a purpose of the exhibition?}
\begin{choices}(2)
  \item To call for environmental protection.
  \item To celebrate the cultural unity of China.
  \item To explore the daily life of ethnic groups.
  \item To preserve the materials and techniques.
\end{choices}
\question{22. What will happen if you cancel your reservation 3 hours in advance?}
\begin{choices}(2)
  \item Your money will be returned.
  \item Your account will not be affected.
  \item Your appointment will be rescheduled.
  \item Your cancellation will not be accepted.
\end{choices}
\question{23. What are the visitors required to do when entering the museum?}
\begin{choices}(2)
  \item Take personal carers.
  \item Pay with their QR code.
  \item Pass a security check.
  \item Confirm their reservation.
\end{choices}
\passagelabel{B}
\par When thinking of TCM (Traditional Chinese Medicine) doctors, one might picture an elderly man with a long beard sitting in front of a medicine chest. However, He Suye, a young and talented physician in the latest medical-romance drama The Best Thing, challenges that traditional image. Played by the 27-year-old actor Zhang Linghe, Mr. He is a professional TCM specialist with a gentle personality, but outside of work, he enjoys hobbies like skateboarding and malatang, a type of hotpot, reshaping common views of TCM doctors.\par
\par The Best Thing follows Shen Xifan, a hotel manager suffering from frequent headaches and trouble sleeping. In her search for a cure, she meets Mr. He, and their relationship deepens through medical consultations. To play a genuine TCM practitioner, Zhang studied Chinese medicine with practicing physicians, learning traditional diagnostic(诊断的) and treatment methods. This process helped him gain a deeper appreciation for TCM's integrated philosophy, which emphasizes the interconnectedness of the body, mind, and environment, viewing health as a balance between these elements.\par
\par The drama creatively integrates TCM knowledge by naming each episode after a type of Chinese herbal medicine, connecting TCM's healing methods to the storyline. For example, in one episode, Mr. He performs "wet cupping" to treat a breathing emergency, showcasing TCM's role in critical care.\par
\par The drama particularly appeals to younger viewers. With the growing integration of TCM and Western medicine, as well as expanding international exchanges, TCM practitioners and students now have more diverse career opportunities. For instance, acupuncture (针灸) is now covered in health insurance plans in many countries. Zhang Linghe believes that passion is the driving force behind a successful career-just like his character, He, who remains committed to his path in TCM despite numerous challenges, including initial opposition from his father. Zhang hopes to introduce TCM to a global audience, especially younger generations, through the drama's release on international platforms.\par
\question{24. What do we know about He Suye?}
\begin{choices}(2)
  \item He challenges the principles of TCM.
  \item He is an actor starring in The Best Thing.
  \item He presents a modern image of a TCM doctor.
  \item He has the same hobbies as other TCM doctors
\end{choices}
\question{25. Why did Zhang Linghe learn TCM techniques?}
\begin{choices}(2)
  \item To prepare for the role.
  \item To treat his own illnesses.
  \item To become a practicing physician.
  \item To appreciate TCM's philosophy.
\end{choices}
\question{26. What is the author's purpose in mentioning acupuncture?}
\begin{choices}(2)
  \item To explain the diverse changes in future careers.
  \item To demonstrate TCM's international recognition.
  \item To show the drama's popularity among younger viewers.
  \item To emphasize TCM's advantages over Western medicine.
\end{choices}
\question{27. What can be a suitable title for the text?}
\begin{choices}(2)
  \item A Traditional Approach to TCM
  \item A New Vision for TCM's Future
  \item Exploring the Ancient Roots of TCM
  \item Redefining TCM for a Global Audience
\end{choices}
\passagelabel{C}
\par Among world languages, English places perhaps the highest value on the expressions of thanks, especially in formal and business settings. So it's logical that English speakers tend to conflate the expression of thankfulness with the feeling of thankfulness. If somebody didn't say "thanks", we often think, they must not have been very grateful. And yet, in English, "thank you" gets used so frequently that it can easily start to feel habitual.\par
\par But do we say it as much as we think in our casual interactions with family and friends? A new research paper titled Universals and Cultural Diversity in the Expression of Thankfulness suggests that no, we do not.\par
\par The research team behind the paper looked for expressions of thankfulness in daily informal interactions across eight languages: English, Italian, Polish, Russian, Lao, Cha'palaa, Murrinhpatha, and Siwu. The findings suggest the universal way in which thankfulness is expressed is not through giving thanks, but rather through participation in a social agreement in which mutual (相互的)assistance and working together are the norm. In fact, the practice of giving thanks seems to have little impact on the practice of mutual exchange.\par
\par The practice of giving thanks also varies culturally. In India, for example, people — especially when they are your elders, relatives, or close friends — tend to feel that by thanking them, you're weakening your closeness with them and creating formality and distance that shouldn't exist. In contrast, in England, expressing thankfulness through words is seen as a positive strengthening of relationships, promoting goodwill and connection.\par
\par Nick Enfield, lead researcher of the study, suggests that most of the 6,000 to 7,000 world languages don't have a word for "thank you". "When we consider the negative effects of thanking close relations alongside the fact that many languages are spoken in very small communities where speakers are likely to know each other, that claim starts to feel quite reasonable."\par
\question{28. What does the underlined word "conflate" in paragraph 1 mean?}
\begin{choices}(2)
  \item Compare.
  \item Replace.
  \item Figure out.
  \item Mix up.
\end{choices}
\question{29. How do people express thankfulness universally according to the study?}
\begin{choices}(2)
  \item By saying thanks directly.
  \item By making social agreements.
  \item By using body language.
  \item By helping and cooperating.
\end{choices}
\question{30. How do Indians view the practice of giving thanks to relatives?}
\begin{choices}(2)
  \item It weakens their bond.
  \item It reduces their social distance.
  \item It enhances their connection.
  \item It promotes their mutual trust.
\end{choices}
\question{31. What does Nick Enfield imply in the last paragraph?}
\begin{choices}(2)
  \item Many cultures do not emphasize thankfulness.
  \item Thankfulness might not be expressed in words.
  \item Small communities prefer to thank close friends.
  \item The negative effects of thanking should be avoided.
\end{choices}
\passagelabel{D}
\par Precision agriculture is one of the most significant trends in modern farming. One particularly promising application involves the use of real-time imaging to visualize how tiny drops of liquid are forming on plant surfaces. Unlike conventional agriculture where farmers often rely on general estimates and uniform application of pesticides (杀虫剂) across the entire field, this technology helps farmers maximize productivity while minimizing inputs.\par
\par This development originated from the work of Vishnu Jayaprakash, who had some family exposure to the process of spraying (喷洒) crops and recognized that the main challenge in that process is how to make water-based sprays interact effectively with surfaces of plants. Several factors influence this process, including the pressure, flow and nozzle (喷嘴) design of the application device, and the temperature, sunlight intensity, wind and other conditions at the time of spraying.\par
\par To effectively address this issue, Jayaprakash developed a two-camera system that can be connected to typical spraying equipment and use imaging of the sprayed solution on the target plants to determine necessary adjustments for the best degree of coverage with appropriately sized tiny drops. The system can be used to make on-the-go adjustments in order to get the spray coverage "just right". This allows farmers to use only 20-30\% of the product per acre (英亩) while maintaining key functions like pest control, fertilization, and growth regulation.\par
\par The advancement in precision agriculture is helping farmers reduce costs while maintaining crop protection. Susan Scheufele, an extension researcher at the University of Massachusetts, has conducted field trials with this technology on vegetable crops. In trials for cotton leaf removal, it has shown positive results even when used at 50\% of the minimum recommended application rate on the product label.\par
\par The primary benefit of precision agriculture methods lies beyond cost savings. Despite the fact that this technology still has a long way to go, its adoption ensures that chemical remains in crops are well below harmful levels, safeguarding ecosystems and food safety.\par
\question{32. What can be learned about precision agriculture in paragraph 1?}
\begin{choices}(2)
  \item It visualizes plants more clearly.
  \item It uses pesticides more accurately.
  \item It continuously monitors plant growth.
  \item It estimates the need of the entire field.
\end{choices}
\question{33. How does Jayaprakash's two-camera system work?}
\begin{choices}(2)
  \item By guiding spraying equipment to target plant surfaces.
  \item By calculating the precise amount of pesticide required.
  \item By making real-time adjustments to enhance spray coverage.
  \item By using imaging to closely monitor plant growth and health.
\end{choices}
\question{34. Why does the author mention Susan Scheufele's field trials?}
\begin{choices}(2)
  \item To analyze the responses of different crops.
  \item To compare the application rates across crops.
  \item To illustrate the effectiveness of the technology.
  \item To highlight the limitations of traditional farming.
\end{choices}
\question{35. What is the author's attitude toward precision agriculture?}
\begin{choices}(2)
  \item Intolerant.
  \item Doubtful.
  \item Unclear.
  \item Approving.
\end{choices}
\sectitle{第二节（共5小题；每小题2.5分，满分12.5分）}
\par 根据短文内容，从短文后的选项中选出能填入空白处的最佳选项。选项中有两项为多余选项。\par
\par The teenage years are marked by intense changes — physical, social, emotional, and developmental — that can be difficult to deal with at times. \blanknum{36} As you grow older, conflicts with family often increase, mainly due to the desire for more independence. Peer (同龄人) pressure also plays a major role, as teenagers are more likely to be influenced by friends while testing boundaries and seeking autonomy.\par
\par Conflicts during this time can also arise from the intensity of your emotions and your evolving ability to think abstractly. As you work on forming your individual identity, you may become more focused on the world around you, trying to make sense of your own feelings and form opinions on various issues. \blanknum{37} As your outlook changes, these struggles may make it more difficult to navigate relationships with family and cope with peer pressure.\par
\par While conflicts are a natural part of this stage of life, there are ways to manage and resolve them productively \blanknum{38} By expressing your thoughts and feelings clearly, you can prevent misunderstandings and build mutual respect. Equally important is active listening — understanding others' outlook can help you find common ground and reach a resolution.\par
\par \blanknum{39} Staying calm allows you to think more clearly and avoid worsening emotions. Techniques such as deep breathing or temporarily stepping away from the situation can help manage intense feelings. Additionally, turning to trusted adults or peers who have faced similar challenges can help you gain practical strategies for managing emotions\par
\par While conflicts are a normal part of growing up, they can also be opportunities for personal development and self-discovery. \blanknum{40}\par
\begin{choices}(1)
  \item Emotional regulation is also essential.
  \item The most effective strategy is communication.
  \item These conflicts are caused by misunderstandings.
  \item With so much going on, it's no surprise conflicts often arise.
  \item This process can lead to internal and external struggles as you grow.
  \item Take it positively, and you can resolve conflicts and boost your growth.
  \item Recognizing the signs helps you become more aware next time a conflict emerges.
\end{choices}
\parttitle{第三部分 语言知识运用（共两节，满分 30 分）}
\sectitle{第一节（共15小题；每小题1分，满分15分）}
\par 阅读下面短文，从短文后各题所给的 A、B、C 和 D 四个选项中，选出可以填入空白处的最佳选项，\par
\par My early school years were a struggle. I had little interest in many subjects and couldn't \blanknum{41} what I was capable of. Since I had been unsuccessful in almost all I attempted, I had lost my \blanknum{42}. There was a time when I would never \blanknum{43} with anyone around me.\par
\par Things didn't \blanknum{44} until one sunny afternoon when I saw kids swimming and laughing in a community pool. At that moment, I suddenly felt a sense of \blanknum{45} and realized that I needed to make a change and \blanknum{46} myself. I thought it might be helpful to learn a new \blanknum{47} and be part of a different group of people. So, I \blanknum{48} a swimming course.\par
\par In the first few weeks, my instructor was quite \blanknum{49} my sincerity and abilities. This was a vital moment — it made me realize how \blanknum{50} it is for a person to discover their \blanknum{51} and fully pursue it. As the days went by, I began swimming \blanknum{52} in local swim meets and even won a few medals. Surprisingly, I made\par
\par \blanknum{53} improvements as well, and finally went to college. When my parents expressed their complete admiration for my performance, my dignity was largely \blanknum{54}.\par
\par Now, I work as a lifeguard, a mountaineer, and a coach. If there's one thing I've learned, it's that you just need to give yourself a \blanknum{55}. Everything I've done has been driven by a single goal-to realize my potential.\par
\cz{41}{deny}{recognize}{accept}{abandon}
\cz{42}{patience}{focus}{dignity}{sympathy}
\cz{43}{fight}{negotiate}{compete}{interact}
\cz{44}{settle}{work}{improve}{worsen}
\cz{45}{regret}{peace}{achievement}{inspiration}
\cz{46}{reinvent}{express}{surprise}{question}
\cz{47}{concept}{routine}{skill}{insight}
\cz{48}{put together for}{kept up with}{dropped out of}{signed up for}
\cz{49}{impressed with}{concerned with}{dismissive of}{critical of}
\cz{50}{alarming}{crucial}{unforgettable}{confusing}
\cz{51}{potential}{success}{secret}{fortune}
\cz{52}{nervously}{competitively}{properly}{smoothly}
\cz{53}{administrative}{technological}{academic}{economic}
\cz{54}{restored}{ruined}{observed}{revealed}
\cz{55}{blessing}{praise}{reward}{chance}
\sectitle{第二节（共10小题；每小题1.5分，满分15分）}
\par 阅读下面短文，在空白处填入 1 个适当的单词或括号内单词的正确形式。\par
\par A standout tradition in Guangzhou is Yum Cha, \blanknum{56} (literal) meaning "drinking tea". Yum Cha is more than just a drink — it's a kind of social activity and a distinct characteristic of Lingnan culture, \blanknum{57} (offer) in family gatherings, business meetings, or casual chats among friends.\par
\par Early-rising \blanknum{58}(local), especially seniors, usually greet each other with "Have you drunk tea?" to start the day. A native saying, "The only thing \blanknum{59} can interrupt a Yum Cha gathering is a typhoon," jokingly but vividly \blanknum{60}(show) the strong commitment of Guangzhou people to their daily tea tradition. In teahouses or restaurants, people come together \blanknum{61}(explore) a wide variety of traditional bite-sized foods, enjoying the social atmosphere and rich cultural heritage (遗产) that accompany the meal, which adds \blanknum{62} touch of depth and enjoyment to the tea gathering.\par
\par Over the last ten years, the tradition of Yum Cha \blanknum{63} (accept) by many young people as a way to reconnect with their roots. Moreover, they are innovating on the classic Yum Cha by experimenting \blanknum{64} low-sugar, low-oil, and high-fiber dim sum (点心), using recyclable and eco-friendly tea ware. This \blanknum{65} (combine) of tradition and innovation helps ensure that Yum Cha culture remains relevant in today's fast-paced world.\par
\parttitle{第四部分 写作（共两节，满分 40 分）}
\sectitle{第一节（满分15分）}
\par 假定你是李华，上周你参加了学校举办的以“绿色能源与未来生活”为主题的科学大讲堂活动。请你给外教 Ryan 写一封电子邮件，分享这次经历，内容包括：\par
\par (1) 活动介绍;\par
\par (2) 你的收获。\par
\par 注意:\par
\par (1) 写作词数应为 80 左右；(2) 请按如下格式在答题纸的相应位置作答，\par
\par\vspace{0.4em}
\begin{center}\begin{tabular}{{|l|}}
\hline
Dear Ryan, \\
Yours, Li Hua \\
\hline
\end{tabular}\end{center}
\vspace{0.4em}
\sectitle{第二节（满分25分）}
\par 阅读下面材料，根据其内容和所给段落开头语续写两段，使之构成一篇完整的短文，\par
\par On a bit erly cold January morning in Alexandria, Minnesota, the temperatures had dropped dangerously below freezing point, and biting winds made conditions even worse. Overnight, a thick blanket of snow had quietly covered the town, streets were deserted, and a sharp silence hung in the air. People stayed indoors, enjoying hot drinks and keeping their fireplaces burning brightly, cautious of the bitter cold that threatened frostbite and hypothermia(失温).\par
\par The town was quiet as people remained indoors. However, Eleanor, an 83-year-old lady who lived alone, worried the birds would go hungry in this freezing weather. As she struggled through deep snow, her boots had suddenly slipped on hidden ice beneath, causing her to fall. Unable to rise with wet gloves and numb (麻木的) fingers, she had been stuck in the snow. Unfortunately, she could turn to nobody for help.\par
\par Originally, Eric Nelson had no intention of stepping outside that morning. He stood at his kitchen window, briefly observing the frozen landscape and rubbing (搓) his hands together to stay warm. However, as he thought about the possibility of frozen water pipes—an inconvenience that could make the already unbearable cold worse—he decided to step out quickly and inspect them.\par
\par As Eric opened his front door, an icy wind struck his face, nearly taking his breath away. To his surprise, he noticed his elderly neighbor Eleanor's small dog, Maya, was already standing anxiously on his doorstep, shaking slightly but barking insistently. She moved her feet anxiously as her breath formed small white clouds in the cold air. Snow stuck to her fur, suggesting she had been outside for longer than was safe in such extreme temperatures.\par
\par Eric instantly sensed something was seriously wrong. Eleanor rarely went outside, especially in such extreme weather. Feeling increasingly uneasy, he quickly followed the anxious dog as she dashed urgently toward Eleanor's backyard. When he arrived, his heart sank. Eleanor lay helplessly on the snow-covered ground near a bird feeder, pale-faced and shaking uncontrollably. 注意：\par
\par （1）续写词数应为 150 左右；\par
\par (2) 请按如下格式在答题卡的相应位置作答,\par
\writingpara{Paragraph 1: Without hesitation, Eric rushed to Eleanor: \blanknum{}}
\writinglines{8}
\writingpara{Paragraph 2: Two days later; Eleanor knocked on Eric's door after being released from the hospital.}
\writinglines{8}
\end{document}
